% ================================================================
%  SESIÓN 07 — Perímetro de zonas urbanas
%  Almería en Cifras y Formas · Matemáticas 1.º ESO
%  IES Al-Ándalus · 2025-2026
% ================================================================
\documentclass[aspectratio=169, 12pt, spanish]{beamer}
\usepackage{almeria-beamer}

\renewcommand{\SessionNumber}{07}
\renewcommand{\SessionTitle}{Perímetro de zonas urbanas}
\renewcommand{\SessionName}{Sesión 07}
\renewcommand{\SessionObjective}{%
  Calcular el perímetro de diferentes figuras geométricas y aplicarlo
  a zonas urbanas de Almería modelizadas como polígonos.%
}
\renewcommand{\BlockName}{Bloque 2: Geometría urbana}

\begin{document}

% ---------------------------------------------------------------
%  PORTADA
% ---------------------------------------------------------------
\begin{frame}[plain]
  \titlepage
\end{frame}

% ---------------------------------------------------------------
%  ÍNDICE
% ---------------------------------------------------------------
\begin{frame}{Contenidos de hoy}
  \begin{columns}[t]
    \begin{column}{0.48\textwidth}
      \begin{itemize}
        \item Definición de perímetro
        \item Analogía: la valla del parque
        \item Fórmulas de perímetro
        \item Pasos para calcular
        \item Unidades de longitud
      \end{itemize}
    \end{column}
    \begin{column}{0.48\textwidth}
      \begin{itemize}
        \item Ejemplos en Almería
        \item Problema paso a paso
        \item Ejercicios Bloom
        \item Error frecuente
        \item Evidencia del día
      \end{itemize}
    \end{column}
  \end{columns}
  \vspace{0.4cm}
  \begin{notabox}[Niveles de Bloom]
    \bloomtag{Recordar}\quad
    \bloomtag{Comprender}\quad
    \bloomtag{Aplicar}
  \end{notabox}
\end{frame}

% ---------------------------------------------------------------
\seccionframe{¿Qué es el perímetro?}
% ---------------------------------------------------------------

% ---------------------------------------------------------------
%  FRAME 3 — Definición
% ---------------------------------------------------------------
\begin{frame}{Perímetro: el borde de la figura}
  \begin{columns}[c]
    \begin{column}{0.54\textwidth}
      \begin{definicionbox}{Perímetro}
        El \textbf{perímetro} de una figura plana es la
        \textbf{suma de las longitudes de todos sus lados}.
        Mide el \emph{contorno}, no el interior.
      \end{definicionbox}

      \vspace{0.4cm}
      \begin{notabox}[Analogía]
        \small El perímetro es como la \textbf{valla} que rodea un parque:
        mides cuánta valla necesitas, no cuánto césped hay dentro.
      \end{notabox}

      \vspace{0.3cm}
      \textbf{Unidades:} m, km, cm, mm \ldots
      \\
      \textbf{Nunca:} m$^2$, km$^2$ (eso sería área).
    \end{column}
    \begin{column}{0.42\textwidth}
      \begin{center}
      \begin{tikzpicture}[scale=0.80]
        % Parque (rectángulo)
        \draw[zona] (0,0) rectangle (3.2,2.0);
        % Valla (borde grueso destacado)
        \draw[AlmTerra, line width=3pt] (0,0) rectangle (3.2,2.0);
        % Césped interior
        \fill[AlmVerde!20] (0.1,0.1) rectangle (3.1,1.9);
        \node[font=\small\bfseries, color=AlmVerde] at (1.6,1.0) {ÁREA};
        \node[font=\small\bfseries, color=AlmTerra] at (1.6,-0.35) {PERÍMETRO (valla)};
        % Flechas de las dimensiones
        \draw[acota] (0,-0.6) -- (3.2,-0.6)
          node[midway, below, font=\scriptsize] {$l$};
        \draw[acota] (3.6,0) -- (3.6,2.0)
          node[midway, right, font=\scriptsize] {$w$};
      \end{tikzpicture}
      \end{center}
    \end{column}
  \end{columns}
\end{frame}

% ---------------------------------------------------------------
%  FRAME 4 — Fórmulas
% ---------------------------------------------------------------
\begin{frame}{Fórmulas de perímetro}
  \begin{columns}[t]
    \begin{column}{0.50\textwidth}
      \formuladest[AlmAzulM]{P_{\text{rect}} = 2l + 2w = 2(l+w)}
      \vspace{0.1cm}
      \formuladest[AlmVerde]{P_{\text{cuad}} = 4l}
      \vspace{0.1cm}
      \formuladest[AlmTerra]{P_{\text{triáng}} = a + b + c}
    \end{column}
    \begin{column}{0.46\textwidth}
      \formuladest[BAnalizar]{P_{\text{reg.}} = n \cdot s}
      \vspace{0.1cm}
      \begin{teoriabox}
        \small
        Donde $n$ = número de lados y $s$ = longitud de cada lado.\\[6pt]
        Para cualquier polígono:
        \[P = \sum_i l_i\]
        (suma de todos los lados).
      \end{teoriabox}
    \end{column}
  \end{columns}
\end{frame}

% ---------------------------------------------------------------
%  FRAME 5 — Pasos para calcular
% ---------------------------------------------------------------
\begin{frame}{Cómo calcular el perímetro: paso a paso}
  \begin{columns}[t]
    \begin{column}{0.50\textwidth}
      \begin{pasobox}[1]
        Identifica \textbf{todos los lados exteriores}
        de la figura. No olvides ninguno.
      \end{pasobox}
      \vspace{0.3cm}
      \begin{pasobox}[2]
        Asegúrate de que \textbf{todas las medidas}
        están en la \textbf{misma unidad}.
        Si hay metros y centímetros, convierte primero.
      \end{pasobox}
    \end{column}
    \begin{column}{0.46\textwidth}
      \begin{pasobox}[3]
        \textbf{Suma} todas las longitudes.
      \end{pasobox}
      \vspace{0.3cm}
      \begin{pasobox}[4]
        Escribe el resultado con su \textbf{unidad}
        (m, km, cm…).
        El perímetro es siempre una \emph{longitud}.
      \end{pasobox}
      \vspace{0.3cm}
      \begin{notabox}[Nunca olvides]
        \small La unidad del perímetro es de
        \textbf{una sola dimensión}: m, no m².
      \end{notabox}
    \end{column}
  \end{columns}
\end{frame}

% ---------------------------------------------------------------
\seccionframe[BAplicar]{Perímetros en Almería}
% ---------------------------------------------------------------

% ---------------------------------------------------------------
%  FRAME 6 — Ejemplos en Almería
% ---------------------------------------------------------------
\begin{frame}{Ejemplos en Almería}
  \begin{columns}[t]
    \begin{column}{0.50\textwidth}
      \begin{ejemplobox}
        \small
        \textbf{Parque Nicolás Salmerón}
        (rectángulo aprox. 400 m $\times$ 80 m):
        \[P = 2(400 + 80) = 2 \times 480 = \mathbf{960\ m}\]

        \vspace{0.15cm}
        \textbf{Plaza triangular hipotética}
        (lados 120 m, 85 m, 100 m):
        \[P = 120 + 85 + 100 = \mathbf{305\ m}\]

        \vspace{0.15cm}
        \textbf{Plaza cuadrada} (lado 60 m):
        \[P = 4 \times 60 = \mathbf{240\ m}\]
      \end{ejemplobox}
    \end{column}
    \begin{column}{0.46\textwidth}
      \begin{center}
      \begin{tikzpicture}[scale=0.70]
        % Rectángulo (parque)
        \draw[zona] (0,0) rectangle (4.0,0.8);
        \node[above, font=\tiny\bfseries, color=AlmAzulM] at (2.0,0.8) {Parque};
        \draw[acota] (0,-0.35) -- (4.0,-0.35)
          node[midway,below,font=\tiny]{400 m};
        \draw[acota] (4.3,0) -- (4.3,0.8)
          node[midway,right,font=\tiny]{80 m};
        % Triángulo
        \draw[zonat] (0,1.5) -- (2.2,1.5) -- (1.5,2.5) -- cycle;
        \node[above, font=\tiny\bfseries, color=AlmTerra] at (1.1,2.5) {Plaza $\blacktriangle$};
        \node[below, font=\tiny, color=AlmGris] at (1.1,1.45) {120 m};
        % Cuadrado
        \draw[rectov, fill=AlmVerde!15] (2.8,1.5) rectangle (4.0,2.7);
        \node[above, font=\tiny\bfseries, color=AlmVerde] at (3.4,2.7) {Plaza $\square$};
        \node[below, font=\tiny, color=AlmGris] at (3.4,1.45) {60 m};
      \end{tikzpicture}
      \end{center}
    \end{column}
  \end{columns}
\end{frame}

% ---------------------------------------------------------------
%  FRAME 7 — Problema completo paso a paso
% ---------------------------------------------------------------
\begin{frame}{Problema resuelto: manzana de Almería}
  \begin{definicionbox}{Enunciado}
    \small Una manzana rectangular del centro de Almería mide
    \textbf{180 m} de largo y \textbf{120 m} de ancho.
    Calcula su perímetro y el tiempo necesario para darle
    la vuelta caminando a 5 km/h.
  \end{definicionbox}

  \vspace{0.2cm}
  \begin{columns}[t]
    \begin{column}{0.48\textwidth}
      \begin{pasobox}[1]
        \small Identificar la fórmula:
        \[P = 2(l + w)\]
      \end{pasobox}
      \vspace{0.2cm}
      \begin{pasobox}[2]
        \small Sustituir valores:
        \[P = 2(180 + 120) = 2 \times 300 = \mathbf{600\ m}\]
      \end{pasobox}
    \end{column}
    \begin{column}{0.48\textwidth}
      \begin{pasobox}[3]
        \small Convertir a km para calcular el tiempo:
        \[600\ \text{m} = 0{,}6\ \text{km}\]
      \end{pasobox}
      \vspace{0.2cm}
      \begin{pasobox}[4]
        \small Calcular el tiempo:
        \[t = \frac{d}{v} = \frac{0{,}6\ \text{km}}{5\ \text{km/h}}
          = 0{,}12\ \text{h} = \mathbf{7{,}2\ \text{min}}\]
      \end{pasobox}
    \end{column}
  \end{columns}
\end{frame}

% ---------------------------------------------------------------
%  FRAME 8 — Comparativa de perímetros (PGFPlots)
% ---------------------------------------------------------------
\begin{frame}{Comparativa: perímetros de zonas de Almería}
  \begin{center}
  \begin{tikzpicture}
    \begin{axis}[
      almeria bar,
      height=5.5cm, width=10cm,
      ymin=0, ymax=1100,
      ylabel={Perímetro (m)},
      symbolic x coords={Parque Salmerón, Plaza cuadrada, Plaza triangular},
      xtick=data,
      xticklabel style={font=\scriptsize, align=center, text width=2.5cm},
      ytick={0,200,400,600,800,1000},
      bar width=1.2cm,
    ]
      \addplot[fill=AlmAzulM, draw=AlmAzul] coordinates {
        (Parque Salmerón, 960)
        (Plaza cuadrada, 240)
        (Plaza triangular, 305)
      };
    \end{axis}
  \end{tikzpicture}
  \end{center}
  \vspace{-0.2cm}
  \begin{notabox}[Observa]
    \small El Parque Nicolás Salmerón tiene el mayor perímetro (960 m)
    porque sus dimensiones totales son mucho mayores, aunque su forma
    sea rectangular.
  \end{notabox}
\end{frame}

% ---------------------------------------------------------------
\seccionframe[BRecordar]{Ejercicios — RECORDAR}
% ---------------------------------------------------------------

% ---------------------------------------------------------------
%  FRAME 9 — Ejercicio RECORDAR
% ---------------------------------------------------------------
\begin{frame}{Ejercicio 1 — RECORDAR \dificultad{1}}
  \begin{recordarbox}
    Calcula el perímetro de cada figura:
    \begin{enumerate}
      \item Rectángulo de \textbf{8 cm} de largo y \textbf{5 cm} de ancho.
      \item Cuadrado de lado \textbf{7 m}.
      \item Triángulo con lados \textbf{6 m}, \textbf{8 m} y \textbf{10 m}.
    \end{enumerate}
  \end{recordarbox}

  \vspace{0.25cm}
  \begin{columns}[t]
    \begin{column}{0.32\textwidth}
      \begin{center}
      \begin{tikzpicture}[scale=0.7]
        \draw[zona] (0,0) rectangle (2.4,1.5);
        \draw[acota] (0,-0.45) -- (2.4,-0.45)
          node[midway,below,font=\tiny]{8 cm};
        \draw[acota] (2.75,0) -- (2.75,1.5)
          node[midway,right,font=\tiny]{5 cm};
      \end{tikzpicture}
      \end{center}
    \end{column}
    \begin{column}{0.32\textwidth}
      \begin{center}
      \begin{tikzpicture}[scale=0.7]
        \draw[rectov, fill=AlmVerde!15] (0,0) rectangle (1.8,1.8);
        \draw[acota] (0,-0.45) -- (1.8,-0.45)
          node[midway,below,font=\tiny]{7 m};
      \end{tikzpicture}
      \end{center}
    \end{column}
    \begin{column}{0.32\textwidth}
      \begin{center}
      \begin{tikzpicture}[scale=0.7]
        \draw[zonat] (0,0) -- (2.4,0) -- (1.8,1.5) -- cycle;
        \node[below,font=\tiny,color=AlmGris] at (1.2,-0.1) {6 m};
        \node[right,font=\tiny,color=AlmGris] at (2.1,0.75) {8 m};
        \node[left,font=\tiny,color=AlmGris] at (0.9,0.75) {10 m};
      \end{tikzpicture}
      \end{center}
    \end{column}
  \end{columns}
\end{frame}

% ---------------------------------------------------------------
\seccionframe[BComprender]{Ejercicios — COMPRENDER}
% ---------------------------------------------------------------

% ---------------------------------------------------------------
%  FRAME 10 — Ejercicio COMPRENDER
% ---------------------------------------------------------------
\begin{frame}{Ejercicio 2 — COMPRENDER \dificultad{2}}
  \begin{comprenderbox}
    Dos parques de Almería:
    \begin{itemize}
      \item \textbf{Parque A:} rectángulo de \textbf{200 m $\times$ 100 m}.
      \item \textbf{Parque B:} cuadrado de lado \textbf{150 m}.
    \end{itemize}
    \begin{enumerate}
      \item Calcula el perímetro de cada parque.
      \item ¿Cuál tiene mayor perímetro?
      \item ¿Por qué el área \textbf{no} es relevante para saber
            cuánta valla se necesita?
    \end{enumerate}
  \end{comprenderbox}
  \vspace{0.2cm}
  \begin{columns}[c]
    \begin{column}{0.46\textwidth}
      \begin{tikzpicture}[scale=0.7]
        \draw[zona] (0,0) rectangle (3.2,1.6);
        \node[font=\scriptsize\bfseries, color=AlmAzulM] at (1.6,0.8) {Parque A};
        \draw[acota] (0,-0.4) -- (3.2,-0.4)
          node[midway,below,font=\tiny]{200 m};
        \draw[acota] (3.55,0) -- (3.55,1.6)
          node[midway,right,font=\tiny]{100 m};
      \end{tikzpicture}
    \end{column}
    \begin{column}{0.46\textwidth}
      \begin{tikzpicture}[scale=0.7]
        \draw[rectov, fill=AlmVerde!15] (0,0) rectangle (2.4,2.4);
        \node[font=\scriptsize\bfseries, color=AlmVerde] at (1.2,1.2) {Parque B};
        \draw[acota] (0,-0.4) -- (2.4,-0.4)
          node[midway,below,font=\tiny]{150 m};
      \end{tikzpicture}
    \end{column}
  \end{columns}
\end{frame}

% ---------------------------------------------------------------
\seccionframe[BAplicar]{Ejercicios — APLICAR}
% ---------------------------------------------------------------

% ---------------------------------------------------------------
%  FRAME 11 — Ejercicio APLICAR
% ---------------------------------------------------------------
\begin{frame}{Ejercicio 3 — APLICAR \dificultad{3}}
  \begin{aplicarbox}
    \small El \textbf{Parque de las Familias} de Almería tiene forma
    de polígono irregular de 6 lados con las medidas:
    \textbf{250 m, 180 m, 300 m, 200 m, 150 m, 220 m}.

    \begin{enumerate}
      \item Calcula su perímetro total.
      \item Rodear el parque contabiliza como \textbf{1 vuelta}.
            ¿Cuántas vueltas completas caben en \textbf{2 km}?
            ¿Sobra algún metro?
    \end{enumerate}
  \end{aplicarbox}

  \vspace{0.2cm}
  \begin{columns}[c]
    \begin{column}{0.52\textwidth}
      \begin{center}
      \begin{tikzpicture}[scale=0.55]
        \draw[zona]
          (0,1.5) -- (2.5,0) -- (5.5,0.3) --
          (6.5,2.0) -- (5.0,3.5) -- (1.5,3.2) -- cycle;
        \node[font=\tiny, color=AlmGris] at (1.25,0.55) {250 m};
        \node[font=\tiny, color=AlmGris] at (4.0,-0.25) {300 m};
        \node[font=\tiny, color=AlmGris] at (6.25,1.15) {180 m};
        \node[font=\tiny, color=AlmGris] at (5.95,2.85) {200 m};
        \node[font=\tiny, color=AlmGris] at (3.15,3.55) {150 m};
        \node[font=\tiny, color=AlmGris] at (0.45,2.45) {220 m};
        \node[font=\scriptsize\bfseries, color=AlmAzulM] at (3.25,1.75) {Parque};
      \end{tikzpicture}
      \end{center}
    \end{column}
    \begin{column}{0.44\textwidth}
      \begin{pasobox}[1]
        \scriptsize $P = 250 + 180 + 300 + 200 + 150 + 220 = ?$
      \end{pasobox}
      \vspace{0.2cm}
      \begin{pasobox}[2]
        \scriptsize Convertir 2 km a metros: $2\ \text{km} = 2000\ \text{m}$
        \\Vueltas $= 2000 \div P$
      \end{pasobox}
    \end{column}
  \end{columns}
\end{frame}

% ---------------------------------------------------------------
%  FRAME 12 — Error frecuente
% ---------------------------------------------------------------
\begin{frame}{¡Cuidado con este error!}
  \begin{errorbox}
    \begin{itemize}
      \item \incorrecto\
            \textbf{«El perímetro del parque es 960 metros cuadrados»}
            \\[6pt]
            \correcto\
            El perímetro es una \textbf{longitud} (1 dimensión),
            su unidad es \textbf{m, km, cm…}
            \\
            Nunca se mide en m$^2$ ni cm$^2$
            (eso sería \emph{área}).
      \vspace{0.4cm}
      \item \incorrecto\
            \textbf{Sumar lados con distintas unidades sin convertir:}
            \\
            $P = 80\ \text{cm} + 1{,}2\ \text{m} + 0{,}5\ \text{m}$
            sin pasar todo a la misma unidad.
            \\[6pt]
            \correcto\
            Convierte primero: $80\ \text{cm} = 0{,}8\ \text{m}$,
            \\
            $P = 0{,}8 + 1{,}2 + 0{,}5 = \mathbf{2{,}5\ \text{m}}$
    \end{itemize}
  \end{errorbox}
\end{frame}

% ---------------------------------------------------------------
%  FRAME 13 — Resumen
% ---------------------------------------------------------------
\begin{frame}{Resumen de la sesión}
  \begin{resumenbox}
    \begin{itemize}
      \item El \textbf{perímetro} es la suma de todos los lados:
            mide el \emph{contorno}, como una valla.
      \item Fórmulas clave:
        $P_{\text{rect}} = 2(l+w)$, \quad
        $P_{\text{cuad}} = 4l$, \quad
        $P_{\text{triáng}} = a+b+c$, \quad
        $P_{\text{reg}} = n\cdot s$.
      \item 4 pasos: identificar lados → unidad única →
            sumar → escribir con unidad.
      \item Unidad del perímetro: \textbf{longitud} (m, km, cm…),
            \textbf{nunca} m$^2$.
      \item Parque Salmerón: $P \approx 960\ \text{m}$.
    \end{itemize}
  \end{resumenbox}
\end{frame}

% ---------------------------------------------------------------
%  FRAME 14 — Evidencia del día
% ---------------------------------------------------------------
\begin{frame}{Evidencia del día}
  \begin{evidenciabox}
    \textbf{Ficha de Perímetros de Almería}

    \vspace{0.2cm}
    Elige \textbf{3 zonas} de Almería (parques, plazas o manzanas
    del callejero). Para cada una:
    \begin{enumerate}
      \item Dibuja la forma aproximada con las medidas.
      \item Muestra \textbf{todos los pasos} del cálculo
            (no solo el resultado final).
      \item Escribe el resultado con su unidad correcta.
      \item Calcula cuánto tiempo tardarías en rodearla
            caminando a 4 km/h.
    \end{enumerate}
    \vspace{0.2cm}
    \textbf{Criterio de éxito:} los 3 perímetros están bien calculados,
    los pasos son visibles y las unidades son correctas.
  \end{evidenciabox}
\end{frame}

\end{document}
